% Related Work for the SaTML 2027 submission (IEEEtran conference, 10pt).
% Bib keys are restricted to the approved list. No result value is typed here;
% if a number is ever needed it must come from generated/macros.tex.

\section{Related Work}

\paragraph{Randomized smoothing and certified accuracy}
Randomized smoothing certifies an $\ell_2$ radius for the Gaussian-smoothed classifier, which returns the class most probable under input noise \cite{lecuyer2019pixeldp,li2019certified,cohen2019smoothing}. Cohen et al.\ give the radius $\sigma\,\Phi^{-1}(p_L)$ from a lower confidence bound $p_L$ and the \textsc{Certify} procedure, which selects a class from one set of draws and bounds its probability from a second \cite{cohen2019smoothing}. Salman et al.\ train the base classifier against attacks on the smoothed model \cite{salman2019provably}; Yang et al.\ characterize the radii attainable under other norms and noise shapes \cite{yang2020randomized}; Kumar et al.\ certify average class scores and margins \cite{kumar2020certifying}, so the use of confidence or margin information inside smoothing is established. In this literature, certified accuracy counts an input only when the smoothed class equals the label and its radius reaches the target; it does not ask whether the smoothed class agrees with the noise-free prediction of the base model. We keep this estimand and report the raw-clean-anchored variant as a separate operational quantity. That smoothing shrinks decision regions unevenly across classes is known \cite{mohapatra2021hidden}, and abstention-aware certified metrics exist for segmentation \cite{anani2024adaptive}; the certified-wrong rate we report is a reporting recommendation grounded in these observations, not a new metric.

\paragraph{Certified vision-language classification}
Nirala et al.\ certify open-vocabulary CLIP classifiers with randomized smoothing, reuse earlier computation when prompts change, and report that certified accuracy degrades sharply as the noise grows \cite{nirala2024ovc}. Our audit runs in a high-noise regime of the kind they describe, and we do not present that degradation as a finding. Seferis et al.\ certify generative vision-language models by mapping their outputs to an oracle classification task \cite{seferis2025randomized}; we study frozen prototype classifiers. PromptSmooth learns textual prompts \cite{zhou2022coop} under Gaussian noise so that a frozen medical vision-language model certifies better \cite{hussein2024promptsmooth}; it is the direct precedent for any text-side adaptation of a smoothed CLIP classifier, and the shared operators we audit are far more constrained than learned prompts. Carlini et al.\ denoise the input with a diffusion model before a pretrained classifier \cite{carlini2023dds}; this image-side route lies outside our operator family and is not evaluated. In embedding space, Pautov et al.\ smooth normalized embeddings and convert an embedding-boundary distance into an $\ell_2$ radius for prototype classifiers \cite{pautov2022smoothed}, and RetrievalGuard certifies smoothed nearest-neighbor retrieval \cite{wu2022retrievalguard}. Xia et al., in a preprint under review, bound cosine feature deviation by a Gaussian robustness score and convert it into a top-one prototype certificate through the clean margin and prototype correlation \cite{xia2026feature}. None of these feature certificates enters our estimates; every certificate we report is a hard-label Cohen certificate of the smoothed version of its frozen classifier (the supervised controls are not zero-shot).

\paragraph{The modality gap and its corrections}
Liang et al.\ define the modality gap of contrastive encoders such as CLIP \cite{radford2021clip}, attribute it to initialization and contrastive optimization, and show that shifting it changes zero-shot behavior in a task-dependent way \cite{liang2022mind}. Schrodi et al.\ find that a few dimensions drive the gap and that its influence on accuracy is often overshadowed by other factors \cite{schrodi2025twoeffects}, and Fahim et al.\ show that contrastive training separates even samples of one modality \cite{fahim2024contrastive}; Levi and Gilboa show that CLIP's image and text embeddings lie on linearly separable ellipsoid shells not centered at the origin \cite{levi2025double}. That the global gap magnitude does not explain accuracy is therefore established. Post-hoc corrections form the operator families we audit. Text-side mean-centering and the two-sided variant that also recenters image features are normalized zero-shot adaptations of the GR-CLIP calibration \cite{li2025closing}. Chowers et al.\ treat the gap as a robustness feature: under a centered-orthogonality assumption and additive noise in embedding space, translating one modality toward the other provably improves a nearest-neighbor stability measure, and a projected translation preserves clean nearest neighbors \cite{chowers2026bug}. Our audit does not contradict that theorem. It concerns a different object, the hard-label vote of a classifier smoothed with pixel noise, and shows that the normalized zero-shot adaptation of their projected operator cannot change any such vote. Sato et al.\ show that clean gap reduction can change class-relative decisions and induce prediction hubness \cite{sato2026hubness}; the certified-wrong class concentration we report arises without any correction and under Gaussian pixel noise, and must not be read as correction-induced clean hubness. Empirical robustness work aligns image and text representations during robust fine-tuning or by training-free projection \cite{dong2025stabilizing,zhu2025cola,dong2025pathsimplices,dong2025subspaces} or fine-tunes the vision encoder adversarially \cite{schlarmann2024robustclip}. These methods target attack accuracy rather than Cohen certificates and lie outside our frozen shared operator family.

\paragraph{Preregistration and post-selection inference}
Nosek et al.\ argue that preregistration separates hypothesis generation from hypothesis testing \cite{nosek2018preregistration}. Kuchibhotla et al.\ review the bias that arises when the same data select a candidate and estimate its effect, and the simultaneous-inference remedies \cite{kuchibhotla2022postselection}. Our validation split both selected and evaluated candidates, so we report the complete family under one simultaneous band and describe the selected candidate rather than test it. Lakens notes that an equivalence claim needs a prespecified smallest effect of interest and a dedicated procedure \cite{lakens2017equivalence}; the audit registered neither and claims no absence of effect, and the follow-up registered both before its sampling (Section~\ref{sec:followup}).

\paragraph{What is and is not new}
New here are a necessary decision-change condition for a renormalized common text-side translation under hard-label smoothing, stated per draw, with the corollary that the projected-gap operator is decision-invariant; a metric-semantics audit that keeps standard certified accuracy, raw-clean-anchored stability, abstention, and certified-wrong concentration separate and shows that these estimands rank operators differently; and a preregistered audit that reports the complete operator family with simultaneous bands, including the one operator whose band excludes zero at the primary radius, followed by a follow-up, registered with a planning power calculation before its sampling, that delimits that operator's gain on fresh items. Not new, and not claimed, are certified zero-shot or open-vocabulary VLM classifiers and their degradation at high noise \cite{nirala2024ovc}; mean-centering and projected gap corrections \cite{li2025closing,chowers2026bug}; noise-aware prompt adaptation \cite{hussein2024promptsmooth}; finite-sample certification machinery \cite{cohen2019smoothing}; class-specific low-rank adaptation \cite{hu2022lora}, used here only as a supervised positive control; and feature-space or embedding certificates \cite{xia2026feature,pautov2022smoothed}.
